\subsection{Arbeit und Energie}
\Aufgabe{

Ein Elektromotor hebt eine Masse von $m = 10\,$kg um eine Höhe von 5\,m.
    \begin{itemize}
        \item[\bf a)]
        Berechnen Sie die verrichtete Arbeit.
        \item[\bf b)]
        Wieviel Energie wird benötigt, wenn der Wirkungsgrad des Motors $\eta = 80\,\%$ beträgt?
    \end{itemize}
}

\Loesung{
\begin{itemize}
  \item[\bfseries a)]
  \begin{flalign*}
    W &= m \cdot g \cdot h &&\\
    W &= 10\,\mathrm{kg}\cdot 9{,}81\,\mathrm{\frac{m}{s^2}}\cdot 5\,\mathrm{m} &&\\
    W &= 490{,}5\,\mathrm{J} &&
  \end{flalign*}

  \item[\bfseries b)]
  \begin{flalign*}
    \text{Wirkungsgrad} &= \frac{\text{genutzte Energie}}{\text{aufgewendete Energie}} = \eta &&\\
    E_\mathrm{zu} &= \frac{W}{\eta} = \frac{490{,}5\,\mathrm{J}}{0{,}8} = 613{,}125\,\mathrm{J} &&
  \end{flalign*}
\end{itemize}
}

\subsection{Höhenenergie im Stausee}
\Aufgabe{
Ein Stausee enthält $2\,000\,000 \,\mathrm{m^3}$ Wasser auf einer Höhe von $100 \,\mathrm{m}$.\\
 \begin{itemize}
        \item[\bf a)]
Berechne die gespeicherte potenzielle Energie. Gegeben ist die Dichte $\rho = 1000\,\mathrm{kg/m^3}\, \text{und}\, \\ g = 9,81\,\mathrm{m/s^2}$
 \end{itemize}
}
\Loesung{
\begin{itemize}
  \item[\bfseries a)]
  \begin{flalign*}
    E_{\text{pot}} &= m \cdot g \cdot h &&\\
    m &= \rho \cdot V
       = 1000\,\mathrm{\frac{kg}{m^3}} \cdot 2\,000\,000\,\mathrm{m^3}
       = 2\,000\,000\,000\,\mathrm{kg} &&\\
    E_{\text{pot}} &= 2\,000\,000\,000\,\mathrm{kg} \cdot 9{,}81\,\mathrm{\frac{m}{s^2}} \cdot 100\,\mathrm{m} &&\\
    E_{\text{pot}} &= 1{,}962 \cdot 10^{12}\,\mathrm{J} &&
  \end{flalign*}
\end{itemize}
}

\subsection{Bewegungsenergie eines Autos}
\Aufgabe{
Ein Auto mit der Masse 1200\,kg fährt mit 72\,km/h.
\begin{itemize}
        \item[\bf a)]
Berechne die kinetische Energie.
 \end{itemize}
}
\Loesung{
\begin{itemize}
  \item[\bfseries a)]
  \begin{flalign*}
    E_{\text{kin}} &= \frac{1}{2}\, m\, v^2 &&\\
    m &= 1200\,\mathrm{kg} &&\\
    v &= 72\,\mathrm{km/h}
      = 72 \cdot \frac{1000\,\mathrm{m}}{3600\,\mathrm{s}}
      = 20\,\mathrm{m/s} &&\\
    E_{\text{kin}} &= \frac{1}{2}\cdot 1200\,\mathrm{kg}\cdot \bigl(20\,\mathrm{\frac{m}{s}}\bigr)^2 &&\\
    E_{\text{kin}} &= 240\,000\,\mathrm{J} &&
  \end{flalign*}
\end{itemize}
}

\subsection{Elektrische Energie eines Geräts}
\Aufgabe{
Ein Wasserkocher mit 2000\,W läuft für 3\,Minuten.
\begin{itemize}
        \item[\bf a)]
        Berechne, wie viel elektrische Energie er verbraucht hat.
        \item[\bf b)]
        Wie viel Energie ist nötig, um 5\,Liter Wasser von 20\,°C auf 100\,°C zu erwärmen?
    \end{itemize}
}
\Loesung{
    \begin{itemize}
        \item[\bf a)]
        \begin{flalign*}
            P &= 2000\,\mathrm{W} &&\\
            t &= 3\,\mathrm{min} = 180\,\mathrm{s} &&\\
            E &= P \cdot t = 2000\,\mathrm{W} \cdot 180\,\mathrm{s} = 360\,000\,\mathrm{J} &&
        \end{flalign*}

        \item[\bf b)]
        \begin{flalign*}
            m &= 5\,\mathrm{kg} \quad (\text{Dichte Wasser} \approx 1\,\mathrm{kg/L}) &&\\
            c &= 4186\,\mathrm{\frac{J}{kg \cdot K}} \quad (\text{spezifische Wärmekapazität von Wasser}) &&\\
            \Delta T &= 100^\circ\mathrm{C} - 20^\circ\mathrm{C} = 80\,\mathrm{K} &&\\
            Q &= m \cdot c \cdot \Delta T = 5\,\mathrm{kg} \cdot 4186\,\mathrm{\frac{J}{kg \cdot K}} \cdot 80\,\mathrm{K} = 1\,674\,400\,\mathrm{J} &&
        \end{flalign*}
    \end{itemize}
}

\subsection{Energie und Leistung}
\Aufgabe{
Ein Heizgerät hat eine Leistungsaufnahme von $P = 2\,$kW und wird für $t = 3$\,h betrieben.
    \begin{itemize}
        \item [\bf a)] Wie viel Energie wird verbraucht?
        \item [\bf b)] Wenn der Strompreis $p$ bei 0,30\,€/kWh liegt, wie hoch sind die kosten $c$\,?
    \end{itemize}
}
\Loesung{
  \begin{itemize}
    \item[\bf a)]
    \begin{flalign*}
      E &= P\cdot t &&\\
      E &= 2\,\text{kW} \cdot 3\,\text{h} = 6\,\text{kWh} &&
    \end{flalign*}

    \item[\bf b)]
    \begin{flalign*}
      c &= E \cdot p &&\\
      c &= 6\,\mathrm{kWh} \cdot 0{,}30\,\frac{\mathrm{Euro}}{\mathrm{kWh}} = 1{,}80\,\mathrm{Euro} &&
    \end{flalign*}
  \end{itemize}
}

\subsection{Wirkungsgrad}
\Aufgabe{
Ein Generator erzeugt eine elektrische Leistung von $P_{\text{el}}$ = 1\,kW, in dem er eine mechanische Leistung von $P_{\text{mech}}$ = 1,5\,kW aufnimmt.
   \begin{itemize}
     \item [\bf a)]Berechne den Wirkungsgrad
     \item [\bf b)]Wie groß muss $P_{\text{mech}}$ sein, um $P_{\text{el}}$ = 2\,kW zu erzeugen?
   \end{itemize}
}
\Loesung{
  \begin{itemize}
    \item[\bf a)]
    \begin{flalign*}
      \eta &= \frac{P_{\text{el}}}{P_{\text{mech}}}
      = \frac{1\,\text{kW}}{1{,}5\,\text{kW}}
      \approx 0{,}6667
      \Rightarrow \eta \approx 66.67\% &&
    \end{flalign*}

    \item[\bf b)]
    \begin{flalign*}
      P_{\text{mech}} &= \frac{2\,\text{kW}}{0{,}6667} \approx 3\,\text{kW} &&
    \end{flalign*}
  \end{itemize}
}
